\chapter{Wprowadzenie w tematykę pracy}

W ostatnich latach rozwój sensorów i urządzeń treningowych umożliwił sportowcom amatorskim i zawodowym gromadzenie dużych ilości danych telemetrycznych opisujących przebieg aktywności fizycznej. W kolarstwie do najważniejszych zmiennych należą moc generowana przez zawodnika oraz tętno, kadencja i prędkość. Analiza takich danych otwiera nowe możliwości zarówno w zakresie monitorowania obciążeń treningowych, jak i w przewidywaniu poziomu formy sportowej.

Celem niniejszej pracy jest opracowanie modelu predykcyjnego pozwalającego na szacowanie funkcjonalnego progu mocy (Functional Threshold Power, FTP) na podstawie szeregów czasowych rejestrowanych podczas jazdy na rowerze. FTP jest parametrem szeroko stosowanym przez kolarzy do planowania treningów i oceny wydolności, jednak jego tradycyjne wyznaczenie wymaga specjalistycznych testów terenowych lub laboratoryjnych. Proponowane rozwiązanie oparte na uczeniu maszynowym ma na celu umożliwienie estymacji tego parametru bez potrzeby wykonywania dedykowanego testu, a jedynie w oparciu o dane zebrane podczas zwykłych sesji treningowych.

Zakres pracy obejmuje analizę literatury dotyczącej pomiaru i interpretacji FTP, charakterystykę dostępnych źródeł danych telemetrycznych oraz opis metodyki ich przetwarzania. W kolejnych rozdziałach przedstawione zostaną etapy przygotowania zbioru danych, inżynierii cech, budowy modeli regresyjnych oraz oceny ich skuteczności na podstawie wybranych metryk.

Praca kończy się dyskusją uzyskanych wyników, wskazaniem ograniczeń zaproponowanego podejścia oraz propozycjami dalszego rozwoju w kierunku bardziej precyzyjnych i adaptacyjnych algorytmów predykcyjnych.
